\chapter{Fundamente teoretice și studii de specialitate}
\label{cap:preliminarii}

Acest capitol trece în revistă literatura care a stat la baza proiectului: studiile clinice despre terapia bazată pe LEGO și despre interacțiunea copiilor cu autism cu roboții, reperele constructiviste și socioculturale, poziționarea eclectică față de orientările psihologice majore și, în final, conceptele și tehnologiile pe care se sprijină implementarea.

\section{Neurodivergența, suprasolicitarea senzorială și funcțiile executive}
\label{sec:meltdown}

Termenul de \emph{neurodivergență} descrie variațiile naturale ale funcționării neurologice umane — tulburările din spectrul autist (TSA), ADHD, dislexia ș.a. Conform DSM-5 \cite{apa2013}, TSA se caracterizează prin deficite persistente în comunicarea și interacțiunea socială, alături de tipare restrictive și repetitive de comportament. Pentru lucrarea de față, două aspecte sunt esențiale: copiii din spectru procesează cu dificultate indiciile sociale subtile (expresii, ton, contact vizual), ceea ce face interacțiunea umană obositoare, dar răspund foarte bine la medii \emph{structurate și predictibile} — ambele observații conduc direct către robotică. În cazul ADHD, dificultățile centrale țin de funcțiile executive — control inhibitor, memorie de lucru și flexibilitate cognitivă \cite{diamond2013} — funcții pe care Diamond arată că le putem antrena prin exerciții repetate, gradate ca dificultate, exact tipul de activitate pe care un sistem robotic îl administrează cu o consecvență pe care un om nu o poate garanta sesiune după sesiune.

Un fenomen central pentru orice intervenție este criza de suprasolicitare (\emph{meltdown}), distinctă de capriciu: nu este un comportament orientat spre un scop, ci o reacție involuntară de descărcare, declanșată când acumularea de stimuli depășește capacitatea de procesare a copilului. Mazefsky și colaboratorii \cite{mazefsky2013} arată că dificultățile de reglare emoțională sunt un mecanism central al TSA, iar relatările adolescenților autiști culese de Phung și colaboratorii \cite{phung2021} confirmă două concluzii practice: episoadele pot fi adesea \emph{prevenite} dacă semnele timpurii sunt observate la timp, iar reacția cea mai utilă este una calmă, predictibilă și lipsită de judecată. Aceste observații au modelat TriSense pe plan preventiv (structură predictibilă, feedback pozitiv, exercițiu de respirație) și educativ (robotul poate interpreta o variantă stilizată și inofensivă a stării de criză, după care „se liniștește” prin propria rutină de respirație, oferind copilului ocazia de a observa și denumi starea din exterior).

\section{Terapia bazată pe LEGO}

Punctul de plecare conceptual al proiectului este terapia bazată pe LEGO (LBT), formalizată de doctorul Daniel LeGoff pornind de la o observație din sala de așteptare: doi copii cu autism, care nu inițiau interacțiuni sociale, au început spontan să vorbească despre construcțiile lor LEGO \cite{legoff2004}. În studiul inițial pe 47 de copii, urmăriți 24 de săptămâni, LeGoff a măsurat îmbunătățiri pe trei dimensiuni — frecvența autoinițierii contactului social, durata interacțiunilor și reducerea comportamentelor stereotipe — iar un follow-up pe trei ani \cite{legoff2006} a confirmat că efectele se mențin. Independent, grupul lui Baron-Cohen de la Cambridge a comparat LBT cu un alt program de abilități sociale și a găsit ameliorări superioare în grupul LEGO \cite{owens2008}, atribuite caracterului \emph{intrinsec motivant} al materialului: copiii nu „fac terapie”, ci se joacă, iar abilitățile sociale se exersează ca efect secundar. Sistemul LEGO funcționează tocmai pentru că este predictibil — piesele se îmbină într-un singur mod corect, regulile sunt fizice, nu sociale — iar construcția după model implică planificare, memorie de lucru și persistență la sarcină. TriSense preia acest cadru și îl augmentează: robotul devine partenerul de joc care propune modelul, urmărește construcția și oferă feedback. În România, metoda a fost adaptată de psihoterapeuții Marius și Andreea Lumpan \cite{lumpanagerpres2025}, iar experiența mea alături de ei a completat lectura articolelor: am observat direct cum terapeutul calibrează ritmul sesiunii și când intervine înainte de escaladare.

\section{Repere psihologice: constructivism, învățare socioculturală și orientări majore}

Designul activităților se sprijină pe constructivismul lui Piaget și pe teoria socioculturală a lui Vygotsky. Piaget descrie dezvoltarea ca proces activ: copilul își construiește înțelegerea prin acțiune asupra mediului \cite{piaget1969}; manipularea pieselor LEGO este învățare prin acțiune — copilul formulează o ipoteză, o testează, primește feedback și ajustează. Vygotsky mută accentul pe \emph{medierea socială}: învățarea apare în zona proximală de dezvoltare, unde un partener mai competent structurează sarcina și retrage treptat ajutorul \cite{vygotsky1978}. TriSense se plasează ca \emph{mediator} — nu înlocuiește terapeutul, dar oferă scaffolding consecvent: instrucțiuni scurte, același ritual, feedback fără pedeapsă, dificultate crescătoare.

Proiectul este deliberat eclectic față de cele trei mari orientări istorice ale psihologiei. Jocurile care antrenează funcțiile executive (Stop \& Go, construcții gradate, fluență semantică) urmează logica \emph{cognitiv-comportamentală}: practică repetată, feedback imediat, întărire pozitivă \cite{diamond2013}. În schimb, modul de a aborda meltdown-ul urmează o orientare \emph{umanistă}: Rogers descrie relația terapeutică drept cadru de acceptare necondiționată, în care persoana nu este redusă la un set de comportamente corectabile \cite{rogers1951} — de aceea robotul nu „sancționează” eșecul, iar scenariile de criză sunt stilizate, voluntare și urmate de reglare, sub supravegherea unui adult. Această complementaritate este o alegere de design: robotul aduce consistența exercițiilor structurate; terapeutul uman aduce empatia, intuiția și decizia de oprire — el observă semnele de suprasolicitare pe care robotul nu le \emph{simte}. Definirea activităților și a constrângerilor etice a fost discutată cu Marius și Andreea Lumpan \cite{lumpanagerpres2025}, nu ca validare clinică, ci ca verificare că jocurile rămân blânde și utile într-un context ghidat de adult. Un ultim reper vine din metoda \emph{Social Stories} a lui Carol Gray \cite{gray2010}: scenarii sociale scurte, formulate identic la fiecare repetare — principiu pe care TriSense îl aplică în salutul, formulele de încurajare și ritualul de încheiere.

\section{Robotica asistivă în intervențiile pentru copii cu autism}

Dacă dificultățile copiilor cu TSA sunt sociale, de ce am introduce o mașină? Pentru că interacțiunea umană este copleșitoare informațional, în timp ce un robot este \emph{simplu, predictibil și răbdător}. Dautenhahn și Werry, pionierii proiectului AURORA, au argumentat încă din 2004 că tocmai această simplitate controlabilă face din robot un mediator ideal — un pas intermediar între lumea predictibilă a obiectelor și cea impredictibilă a oamenilor \cite{dautenhahn2004}. Sintezele ulterioare confirmă: Scassellati, Admoni și Matarić \cite{scassellati2012} documentează că în prezența roboților copiii cu TSA manifestă contact vizual prelungit, atenție comună, imitație spontană și inițierea interacțiunii, comportamente care se extind și către adulții prezenți; Robins și colaboratorii \cite{robins2005}, cu robotul KASPAR, arată că robotul funcționează ca \emph{mediator social}; iar Kim și colaboratorii \cite{kim2013} găsesc că un copil cu TSA vorbește mai mult cu un adult în prezența unui robot. Sintezele critice impun și prudență — eșantioane mici, design exploratoriu \cite{diehl2012, pennisi2016} — iar Cabibihan și colaboratorii \cite{cabibihan2013} și Ricks și Colton \cite{ricks2010} observă că roboții non-umanoizi sunt adesea mai bine primiți decât cei cu aspect uman realist. Poziționarea TriSense este directă: un robot non-umanoid din piese LEGO, partener de joc și mediator, care administrează activități inspirate din LBT și din probele de funcții executive și produce metrici obiective — fără nicio pretenție de diagnostic.

Mai multe jocuri sunt inspirate din probe consacrate în neuropsihologie, adaptate ludic: „Build the Model” pornește de la \emph{Turnul din Hanoi} (planificare vizuo-spațială), „Stop \& Go” de la paradigma Go/No-Go (control inhibitor), estimarea timpului de la probele de percepție temporală, iar „categoriile rapide” de la fluența semantică. Statutul lor este însă explicit: sunt \emph{exerciții inspirate din procedee standard, nu teste clinice validate} — TriSense nu măsoară inteligența și nu diagnostichează.

\section{Sisteme cyber-fizice și tehnologiile utilizate}

Un \emph{sistem cyber-fizic} (CPS) integrează strâns calculul, comunicația și procesele fizice, cu o componentă de calcul adesea distribuită, în buclă închisă cu mediul \cite{lee2008}. Lee subliniază că provocarea centrală nu este puterea de calcul, ci \emph{orchestrarea}: sincronizarea în timp real a unor componente cu cerințe de fiabilitate pe care software-ul de birou nu le are. TriSense este un CPS în sens deplin — percepe vocea (microfon), construcția (cameră) și apăsările de buton, le procesează printr-o componentă distribuită pe trei dispozitive plus servicii cloud și le transformă în acțiuni fizice (mișcări, animații pe LED-uri, vorbire), caracterul adaptiv venind din stratul de inteligență artificială.

Pe partea hardware, robotul folosește platforma \textbf{LEGO\textsuperscript{®} SPIKE\texttrademark{} Prime} cu firmware-ul open-source \emph{Pybricks} \cite{pybricks}, ales pentru controlul fin al motoarelor și matricei de LED-uri și pentru accesul programabil la portul LPF2. Placa \textbf{LMS-ESP32} \cite{antonsmindstorms} se conectează la acest port și aduce ce îi lipsește hub-ului — Wi-Fi, audio I2S și putere de calcul — rulând MicroPython \cite{micropython}; ea coordonează un amplificator MAX98357A cu difuzor și un microfon INMP441. Pentru viziune, un modul \textbf{ESP32-CAM} de cost redus livrează imagini JPEG la cerere prin HTTP. Comunicația folosește trei mecanisme: \textbf{MQTT} \cite{mqtt} (publish/subscribe, prin brokerul Mosquitto) pentru control, \textbf{LPF2} cu biblioteca \emph{PupRemote} \cite{pupremote} pentru legătura hub--ESP32 (canalele \texttt{cmd} și \texttt{btn}) și conexiuni \textbf{TCP} directe pentru fluxurile audio PCM. Componenta de inteligență artificială se bazează pe modelele \textbf{Google Gemini} \cite{gemini} pentru dialog (cu un \emph{system prompt} care impune un ton cald și propoziții scurte), pentru transcrierea vorbirii și pentru verificarea vizuală a construcțiilor în regim \emph{zero-shot} — modelul primește fotografia și o descriere în limbaj natural a țintei și întoarce un verdict structurat, fără date de antrenare. Sinteza vocală folosește Google Cloud TTS, cu \texttt{edge-tts} ca alternativă locală.
